\section{Loss-indexed selective segmentation}
\label{sec:setup}

\paragraph{Prediction and loss.}
Let \((X,Y)\) be an image and its binary segmentation mask, drawn from an
unknown distribution. A fixed system returns a probability map
\(p(X)\in[0,1]^N\) and a deployed hard prediction \(\Yhat(X)\). For example,
\(\Yhat(X)=\{i:p_i(X)\geq\gamma\}\) for a fixed threshold \(\gamma\), but the
definitions apply to any deterministic deployed decision rule.

Let \(L_x(A,B)\) be the loss incurred when \(B\) is deployed and \(A\) is the
label. We assume that \(L_x\) is measurable and bounded:
\begin{equation}
  0\leq L_x(A,B)\leq D_x<\infty.
  \label{eq:bounded-assumption}
\end{equation}
Boundedness is automatic for \(1-\Dice\) and \(1-\IoU\). A surface loss that
is undefined for an empty mask must be extended before it enters the framework.
The convention is part of \(L\) and must be shared by the confidence computation
and selective-risk evaluation.

\paragraph{Risk and coverage are loss-indexed.}
A confidence \(g:\mathcal X\to\mathbb R\) is oriented so that larger values mean
more confidence. At threshold \(\tau\), define
\begin{align}
  \operatorname{cov}_g(\tau)
    &=\Prb\!\left(g(X)\geq\tau\right), \\
  R_L(\tau;g)
    &=
    \frac{\Ex\!\left[
      L_X(Y,\Yhat(X))\mathbf 1\{g(X)\geq\tau\}
    \right]}
    {\Prb(g(X)\geq\tau)} .
  \label{eq:risk-coverage}
\end{align}
The risk in Equation~\eqref{eq:risk-coverage} is defined only when the
denominator is positive. At a prescribed coverage, randomization within a tied
confidence value attains the requested acceptance probability without imposing
an arbitrary order. Writing \(R_L(c;g)\) for the resulting risk at coverage
\(c\in(0,1]\),
\begin{equation}
  \AURC_L(g)=\int_0^1 R_L(c;g)\,dc .
  \label{eq:aurc-l}
\end{equation}
This is the loss-indexed version of the standard selective-prediction evaluation
\citep{geifman2017selective,selectivenet}. The subscript matters:
\(\AURC_{1-\Dice}\) and \(\AURC_{\HDF}\) evaluate different rankings of harm
and have different units. AURC depends on \(g\) only through its induced ranking,
up to the declared tie rule.

\paragraph{Ideal confidence.}
\begin{definition}[Ideal loss-indexed confidence]
\label{def:ideal}
For a fixed deployed predictor \(\Yhat\), define
\begin{equation}
  \rho_L^\star(x)
  :=\Ex\!\left[L_X(Y,\Yhat(x))\mid X=x\right],
  \qquad
  C_L^\star(x):=-\rho_L^\star(x).
  \label{eq:ideal}
\end{equation}
\end{definition}
If the deployment threshold changes from \(\gamma=0.5\) to \(\gamma=0.3\),
then \(\Yhat\), and therefore \(C_L^\star\), changes. Confidence here is
decision-aware rather than solely a property of an uncertainty distribution.

\begin{proposition}[Optimal selective ordering]
\label{prop:ideal-ranking}
Among all acceptance rules measurable with respect to \(X\), accepting images in
increasing order of \(\rho_L^\star(X)\) minimizes selective risk at every coverage,
up to randomization among ties. Equivalently, ranking by \(C_L^\star\) minimizes
\(\AURC_L\).
\end{proposition}
The result follows by conditioning the accepted loss on \(X\) and exchanging any
accepted point of higher conditional risk with a rejected point of lower conditional
risk. Appendix~\ref{app:theory} gives the measure-theoretic statement and proof.

Definition~\ref{def:ideal} is a target, not an available score. It requires the full
conditional law of \(Y\) given \(x\). Pixel marginals
\(q_i(x)=\Prb(Y_i=1\mid X=x)\) do not identify how pixels vary jointly, and
nonlinear losses such as Dice and HD95 depend on that joint structure. We next make
one working choice explicit.
